% CAPITOLO 4 — TESTING
\chapter{Testing}
\label{cap:testing}

La verifica delle funzionalità e dei requisiti del caso d'uso
\textbf{UC-04 Ordina un Prodotto} (con autenticazione come prerequisito) è
effettuata tramite una suite di test automatizzati \textbf{JUnit 5}
(\emph{Jupiter}) eseguita da \textbf{Maven} (\texttt{mvn verify}) con la
copertura misurata da \textbf{JaCoCo}. La build è riproducibile e la soglia di
copertura è imposta come \emph{quality gate}.

\section{Strumenti e organizzazione}
\label{sec:testing_strumenti}

I test sono scritti con \texttt{JUnit 5} e risiedono in \texttt{src/test/java},
percorso mirror dell'architettura \textbf{BCE} del codice applicativo
(\texttt{src/main/java}). La suite verifica i layer \emph{control}
(\texttt{OrdinaProdottoController}, \texttt{AutenticazioneController}),
\emph{entity} (\texttt{Ordine}, \texttt{Prodotto}, \texttt{BuonoPromozionale}),
i \textbf{pattern} di dominio (Strategy degli sconti, Observer, Singleton,
factory) e la \textbf{persistenza} (\textbf{DAO Demo} e \textbf{FS}). I DAO
\textbf{JDBC} (richiedono un MySQL reale) sono esclusi dalla copertura
automatizzata e validati manualmente, come documentato nel
Capitolo~\ref{cap:persistenza}.

Il plugin \texttt{jacoco-maven-plugin} misura la copertura e impone il
\emph{quality gate}: la build fallisce se la copertura del contatore
\textbf{LINE} scende sotto l'\(70\%\). Gli stessi target esclusi (boundary
JavaFX che richiedono display, bootstrap e DAO JDBC) sono dichiarati nella
configurazione del plugin.

\section{Specifiche dei casi di test}
\label{sec:testing_specifiche}

Di seguito sono descritti i casi di test della suite, che coprono il flusso
completo dell'UC-04.

\subsection{Bean e DTO}
Il test \texttt{BeanTest} verifica la corretta creazione e il ciclo di vita dei
\textbf{Bean} che attraversano i confini BCE (\texttt{OrdineBean},
\texttt{UtenteBean}, \texttt{ProdottoBean}, \texttt{PaymentInfoBean}):
costruttori, valorizzazione iniziale e getter/setter.

\subsection{Controller applicativi}
La suite \texttt{OrdinaProdottoControllerTest} (il test del caso d'uso core)
verifica in stile \emph{We tested that}:
\begin{itemize}
  \item \emph{submit ordine con utente nullo}: rifiutato quando non esiste una
  sessione autenticata;
  \item \emph{submit ordine prodotto non trovato}: un ordine su un prodotto
  inesistente nel catalogo viene scartato;
  \item \emph{submit ordine venditore dal prodotto}: il venditore è risalito
  dal prodotto acquistato (whole-part) e associato correttamente all'ordine;
  \item \emph{riduzione disponibilità}: dopo l'acquisto la scorta del prodotto
  è decrementata di 1 (RF-3);
  \item \emph{quantità eccedente la scorta}: un acquisto oltre le quantità
  disponibili solleva un'eccezione (estensione 3a);
  \item \emph{pagamento autorizzato}: un pagamento approvato dal
  \texttt{PaymentGateway} porta alla creazione dell'ordine in stato
  \texttt{CREATED};
  \item \emph{pagamento rifiutato}: un rifiuto del gateway solleva
  \texttt{BusinessValidationException} senza creare ordini (estensione 6a).
\end{itemize}

\noindent Inoltre:
\begin{description}
    \item[\texttt{AutenticazioneControllerTest}] verifica autenticazione e
    registrazione (credenziali valide/invalide, creazione profilo con il ruolo,
    sessione);
    \item[\texttt{OrdineControllerTest}] verifica il ciclo di vita dell'ordine,
    la transizione di stato (BR-04) e la notifica Observer.
\end{description}

\subsection{Entità e dominio}
\begin{itemize}
    \item \texttt{OrdineTest}: transizione di \texttt{cambiaStato}, applicazione
    del buono, ricalcolo totale e notifica;
    \item \texttt{ProdottoTest}: riduzione della disponibilità e regole di
    quantità;
\end{itemize}

\subsection{Pattern e factory}
\begin{itemize}
    \item \texttt{BuonoPromozionaleStrategyTest} e \texttt{ScontoStrategyFactoryTest}:
    scontistica percentuale/importo fisso e Simple Factory (incluso il ramo di
    errore su tipo ignoto);
    \item \texttt{OrdineLazyFactoryTest}, \texttt{PatternTest}: factory lazy e
    pattern Observer/strategy;
    \item \texttt{ControllerNoArgConstructionTest}: tutti i controller sono
    istanziabili col costruttore no-arg via ServiceLocator (BCE).
\end{itemize}

\subsection{Persistenza}
\begin{itemize}
    \item \texttt{DemoDAOTest} e \texttt{FSDaoTest}: verifica dei DAO Demo
    (in-memory) e Filesystem (JSON);
    \item \texttt{ApplicationModeManagerTest}: risoluzione della modalità
    applicativa (Demo/FS/JDBC) e idempotenza del seed.
\end{itemize}

\section{Tracciabilità Requisiti{-}Test}
\label{sec:testing_tracciabilita}

La matrice collega i Requisiti Funzionali dell'UC-04 (Capitolo~\ref{cap:srs})
alle classi di test JUnit che li verificano.

\begin{table}[H]
    \centering
    \small
    \setlength{\tabcolsep}{4pt}
    \begin{tabular}{@{}llp{4.6cm}@{}} \toprule
        \textbf{Requirement} & \textbf{Passo UC-04} & \textbf{Classe di test}\\ \midrule
        RF-3 / US-3 & selezione e ordine & \texttt{OrdinaProdottoControllerTest}\\
        & passo 3a (scorte) & \texttt{OrdinaProdottoControllerTest}\\
        & passo 6 (pagamento) & \texttt{OrdinaProdottoControllerTest}\\
        & 6a (pagamento negato) & \texttt{OrdinaProdottoControllerTest}\\
        & 4a (buono) & \texttt{BuonoPromozionaleStrategyTest}, \\ & & \texttt{ScontoStrategyFactoryTest}\\
        & BR-04 (stato ordine) & \texttt{OrdineControllerTest}\\
        & Autenticazione & \texttt{AutenticazioneControllerTest}\\
        & Persistenza & \texttt{DemoDAOTest}, \texttt{FSDaoTest}\\
        \bottomrule
    \end{tabular}
    \caption{Matrice di tracciabilità Requisiti{-}Test (JUnit 5/JaCoCo).}
    \label{tab:tracciabilita_test}
\end{table}

\section{Risultati dell'esecuzione}
\label{sec:testing_risultati}

La suite viene eseguita con \texttt{mvn verify}. L'ultima esecuzione (report
JaCoCo, post scope-restrict UC-04) riporta:

\begin{table}[H]
    \centering
    \small
    \begin{tabular}{@{}lr@{}} \toprule
        \textbf{Metrica} & \textbf{Valore} \\ \midrule
        Test eseguiti & 69 \\ \midrule
        Passati & 69 \\ \midrule
        Falliti & 0 \\ \midrule
        Errori & 0 \\ \midrule
        Copertura INSTRUCTION & 73.6\% \\ \midrule
        Copertura LINE & 74.4\% \\ \midrule
        Copertura BRANCH & 57.7\% \\ \midrule
        Copertura METHOD & 74.6\% \\ \midrule
        Copertura CLASS & 96.2\% \\ \bottomrule
    \end{tabular}
    \caption{Risultato complessivo della suite di test (scope UC-04).}
    \label{tab:testriassunto}
\end{table}

Tutti i \textbf{69 test passano} con \textbf{0 errori} e \textbf{0 fallimenti}.
La copertura \textbf{LINE 74.4\%} supera la soglia richiesta (\(70\%\),
imposta nel plugin: la build fallisce in caso contrario). La copertura
\textbf{CLASS 96.2\%} conferma che quasi tutte le classi applicative dello
scope UC-04 sono esercitate.

\section{Copertura per sotto-sistema}
\label{sec:testing_copertura}

La copertura è concentrata sul caso d'uso core: controller
\texttt{OrdinaProdottoController} e \texttt{AutenticazioneController}, entity
\texttt{Ordine}, \texttt{Prodotto}, \texttt{BuonoPromozionale}, pattern
(Strategy/Observer) e persistenza Demo/FS superano la soglia, a garanzia dei
requisiti del UC-04.

\section{SonarCloud e quality gate}
\label{sec:testing_sonar}

Come nelle relazioni di riferimento, è utilizzata l'integrazione
\textbf{SonarCloud} per l'analisi statica (bugs, vulnerabilities, code smells e
duplicati). Il progetto supera il \emph{Quality Gate} configurato ed è
analizzabile con \texttt{mvn sonar:sonar} (plugin \texttt{sonar-maven-plugin}
già presente nel \texttt{pom.xml}).

\begin{quote}
\textbf{Nota di consegna}: da includere lo screenshot del \emph{Quality Gate}
di SonarCloud (stato PASSED) e l'URL pubblico del progetto.
\end{quote}